% -------------------------------------------------------
% Abstrakt / Abstract
% Achtung: Für Anführungszeichen \enquote{} verwenden, nicht "` und "'.
% TODO: Erst nach Fertigstellung der Arbeit final formulieren (max. halbe Seite).

\newcommand{\hsmaabstractde}{\textit{[Entwurf]} Diese Seminararbeit beschreibt Konzeption und Implementierung einer webbasierten Spiellobby, deren Client in React und deren Server in Node.js mit Express und der Bibliothek \texttt{ws} umgesetzt ist. Sämtliche Interaktionen nach dem Login laufen über eine dauerhafte WebSocket-Verbindung: Lobby-Listen und Lobby-Zustände werden in Echtzeit an alle betroffenen Clients gepusht, Freundschaftsanfragen, Einladungen, Online-Status und private Nachrichten erscheinen als Live-Benachrichtigungen, und Lobby-Mitglieder arbeiten über einen gemeinsamen Chat mit Tipp-Anzeige sowie einen geteilten Bereit-Status zusammen. Im Mittelpunkt steht die Frage, wie ein ereignisgetriebener Nachrichtenstrom sauber in das deklarative Rendering-Modell von React integriert werden kann. Die Arbeit zeigt dazu ein Muster aus einem eigenen Hook für den Verbindungslebenszyklus und einem externen Zustand-Store als Brücke zwischen Socket und Benutzeroberfläche.}

% Englisches Abstract wird nicht ausgegeben (nur für die Vorlage definiert)
\newcommand{\hsmaabstracten}{}
